\thispagestyle{empty}
\begin{tikzpicture}[remember picture,overlay]
\begin{scope}[shift={(current page.south west)}]
\fill[otNavy] (0,0) rectangle (210mm,297mm);
\shade[bottom color=otNavyDark,top color=otNavy] (0,0) rectangle (210mm,150mm);
\draw[white,opacity=.05,step=8mm] (0,0) grid (210mm,297mm);
\node[anchor=north west,inner sep=0pt] at (22mm,279mm){\includegraphics[width=60mm]{logo-blanc}};
\node[anchor=west,text=otGreen,font=\otSemi\fontsize{11}{14}\selectfont] at (23mm,239mm){\otCaps[18]{Terminale générale · Option}};
\node[anchor=west,text=white,font=\otBold\fontsize{41}{45}\selectfont] at (21mm,216mm){Mathématiques};
\node[anchor=west,text=white,font=\otLight\fontsize{30}{36}\selectfont] at (22mm,196mm){\NomOption};
\fill[otGreen] (24mm,183mm) rectangle (54mm,184.4mm);
\node[anchor=north west,text=white,text width=156mm,font=\otRegular\fontsize{11}{16}\selectfont] at (24mm,170mm)
{Le cours expliqué\enspace·\enspace Des méthodes résolues\\Des exercices progressifs et leurs corrigés};
\node[anchor=north west,text=white,text width=145mm,font=\otRegular\fontsize{10.5}{16}\selectfont] at (24mm,136mm)
{\otCheck\hspace{3mm}\NbChapitres\ chapitres pour apprendre et réviser\\[3mm]
\otCheck\hspace{3mm}\NbExercices\ exercices, du calcul au problème\\[3mm]
\otCheck\hspace{3mm}Schémas et représentations vectorielles\\[3mm]
\otCheck\hspace{3mm}Dossiers de synthèse et ateliers Python};
\draw[otGreen,line width=2pt] (24mm,51mm)--(60mm,51mm);
\node[anchor=west,text=white,font=\otSemi\fontsize{9}{12}\selectfont] at (24mm,41mm){Programme français · Édition 2026-2027};
\node[anchor=west,text=white,opacity=.7,font=\otLight\fontsize{8}{11}\selectfont] at (24mm,31mm){Manuel indépendant fondé sur les textes du Bulletin officiel};
\node[anchor=west,text=white,font=\otMedium\fontsize{10}{12}\selectfont] at (24mm,16mm){optitab.net};
\end{scope}\end{tikzpicture}\null\clearpage
